\documentclass[12pt]{article}
\usepackage[utf8]{inputenc}
\usepackage[T1]{fontenc}
\usepackage[margin=1in]{geometry}
\usepackage{setspace}
\usepackage{xcolor}
\usepackage[colorlinks=true, urlcolor=blue!60!black, linkcolor=blue!60!black]{hyperref}
\setlength{\parindent}{2em}
\setlength{\parskip}{1.25em}
\renewcommand{\baselinestretch}{1.0}
\begin{document}
\begin{center}
{\LARGE \textit{Griffin v.\ Breckenridge}, 403 U.S.\ 88 (1971)\par}
\vspace{0.8em}
{\large Case Brief --- Four-Page Edition\par}
\end{center}
\vspace{1em}

\section{Memory Jogger}

Under 42 U.S.C. \S~1985(3), Congress constitutionally reached purely private, racially motivated conspiracies without any state action, relying on Section 2 of the Thirteenth Amendment and the right of interstate travel, so long as the conspirators acted with racial or class-based invidiously discriminatory animus (\href{https://supreme.justia.com/cases/federal/us/403/88/}{403 U.S. 88}).

\section{Detailed Case Facts}

On July 2, 1966, Eugene Griffin and other Black residents of Kemper County, Mississippi, were passengers in a car driven by R. G. Grady of Memphis, Tennessee, on highways near DeKalb, Mississippi. Lavon and James Calvin Breckenridge, two white residents of DeKalb, mistakenly believed Grady was a civil rights worker. According to the complaint, they conspired to block the car, detain its occupants, and assault them in order to deprive them of equal protection and equal rights, including free speech, assembly, movement, and use of the highways on equal terms.

The Breckenridges drove their truck into the car's path, forced the occupants out at gunpoint, and threatened to kill them, while James Calvin Breckenridge clubbed each victim about the head, causing serious injuries. Neither defendant acted under color of state law. The victims sued for damages under \S~1985(3), originally \S~2 of the Ku Klux Klan Act of 1871.

\section{Procedural History}

The Southern District of Mississippi dismissed under \textit{Collins v.\ Hardyman}, which had read \S~1985(3) to reach only conspiracies under color of state law. The Fifth Circuit affirmed, voicing ``serious doubts'' about \textit{Collins} but considering itself bound. The Supreme Court reversed and remanded on June 7, 1971.

\section{Judicial Votes}

Reversed and remanded, \textbf{9--0} (8--1 as to Part V-B on the right of interstate travel).
\begin{itemize}
\item \textbf{Majority:} Justice Stewart, joined by Chief Justice Burger and Justices Black, Douglas, Brennan, White, Marshall, and Blackmun, and by Justice Harlan except Part V-B.
\item \textbf{Concurring:} Justice Harlan, joining all but Part V-B.
\item \textbf{Dissenting:} None.
\end{itemize}

\section{Holding}

Section 1985(3) reaches wholly private conspiracies. A plaintiff must allege (1) a conspiracy (2) for the purpose of depriving a person or class of equal protection or equal privileges and immunities, which requires racial, or perhaps otherwise class-based, invidiously discriminatory animus, (3) an act in furtherance, and (4) injury to person or property or deprivation of a citizen's right. Congress had power to reach this conspiracy under Section 2 of the Thirteenth Amendment and its power to protect interstate travel. The Court reserved whether non-racial animus suffices (making ``or perhaps otherwise class-based'' dicta) and whether Section 5 of the Fourteenth Amendment independently supports the statute.

\section{Analysis of Opinions}

\textbf{Majority (Stewart).} Four arguments dismantled \textit{Collins}. The \emph{text} speaks of two or more persons who conspire or go in disguise on the highway, conduct typical of Klan marauders, with no color-of-law requirement. \emph{Precedent} gave Reconstruction statutes a sweep as broad as their language. The \emph{structure} of the 1871 Act addressed state action elsewhere, so a state-action gloss would make this clause redundant. And the \emph{legislative history} showed that Congress meant to reach private conspirators. To keep the statute from becoming a general federal tort law, the Court required invidious, class-based animus, which the racially motivated highway ambush easily satisfied.

On constitutional power, Stewart abandoned the old non-severability rule of \textit{United States v.\ Harris}, asking only whether the statute was valid as applied. Part V-A, relying on the \textit{Civil Rights Cases} and \textit{Jones v.\ Alfred H.\ Mayer Co.}, held that Congress may rationally determine the badges and incidents of slavery and provide a remedy against racially motivated private conspiracies. Part V-B, relying on \textit{United States v.\ Guest} and \textit{Shapiro v.\ Thompson}, held the right of interstate travel protected against private interference; petitioners could prove the attack targeted interstate travel or association with out-of-state civil rights workers. Part V-C reserved the Section 5 question.

\textbf{Harlan (concurring).} Harlan found it unnecessary to rely on the right of interstate travel; the Thirteenth Amendment sufficed.

\section{Examples: Future Applications}

\begin{enumerate}
\item \textbf{Same-Side: Driving Out a Black Family.} White supremacists conspire to vandalize the home of a Black family that moved across state lines, to drive it from town. Private racial animus aimed at basic rights and interstate movement falls within \S~1985(3).
\item \textbf{Same-Side: Attacking a Civil Rights Volunteer.} Private citizens beat an out-of-state white volunteer driving Black residents to a voter rally. As with Grady, the target need not be Black if racial animus drives the conspiracy.
\item \textbf{Opposite-Side: Anti-Union Violence.} Union supporters assault nonunion workers to force union hiring. Under \textit{Carpenters v.\ Scott}, economic animus is not covered, and private interference with First Amendment association is not actionable without state involvement.
\item \textbf{Opposite-Side: Clinic Blockades.} Under \textit{Bray v.\ Alexandria Women's Health Clinic}, opposition to abortion is not class-based animus against women, and incidental effects on interstate travelers do not show a conspiracy aimed at travel.
\item \textbf{Fence-Sitter: Attacking a Religious Minority.} A militia assaults members of a religious minority on an interstate highway to drive their congregation from the state. The conspiracy targets travel and a group, but \textit{Griffin} reserved non-racial animus, and courts have split.
\end{enumerate}

\section{Critique}

\textbf{Originalist and textualist critiques.} Some critics argue that reliance on \textit{Jones} stretches the badges and incidents of slavery beyond the Thirteenth Amendment's original meaning, which they read as limited to chattel slavery and its legal disabilities rather than a general power over private torts. Textualists also note a tension: after insisting on the statute's literal sweep, the Court detached ``equal protection'' from its state-action framework and grafted on an atextual animus requirement.

\textbf{Progressive perspectives.} Civil rights scholars applaud \textit{Griffin} for reviving the Ku Klux Klan Act and grounding protection against private racial terror in the Thirteenth Amendment. Yet critics lament that avoiding Section 5 and inventing an ambiguous animus gatekeeper led to the retrenchment of \textit{Carpenters} and \textit{Bray}, leaving non-racial plaintiffs with little and lower courts confused. The case pits federalism and the state-action boundary against national civil rights enforcement.

\section{Key Quotations}

\begin{quote}
``The language requiring intent to deprive of equal protection, or equal privileges and immunities, means that there must be some racial, or perhaps otherwise class-based, invidiously discriminatory animus behind the conspirators' action. The conspiracy, in other words, must aim at a deprivation of the equal enjoyment of rights secured by the law to all.''
\end{quote}
\noindent --- Justice Stewart, Opinion of the Court, 403 U.S. at 102.

\begin{quote}
``Our cases have firmly established that the right of interstate travel is constitutionally protected, does not necessarily rest on the Fourteenth Amendment, and is assertable against private, as well as governmental, interference.''
\end{quote}
\noindent --- Justice Stewart, Part V-B, 403 U.S. at 105--06.

\begin{quote}
``We can only conclude that Congress was wholly within its powers under \S~2 of the Thirteenth Amendment in creating a statutory cause of action for Negro citizens who have been the victims of conspiratorial, racially discriminatory private action aimed at depriving them of the basic rights that the law secures to all free men.''
\end{quote}
\noindent --- Justice Stewart, Opinion of the Court, 403 U.S. at 105.

\end{document}
